\documentclass[11pt,a4paper]{article}
\usepackage[utf8]{inputenc}
\usepackage[T1]{fontenc}
\usepackage[portuguese]{babel}
\usepackage[a4paper,margin=2.5cm]{geometry}
\usepackage{booktabs}
\usepackage{verbatim}

\title{Segurança de Dados\\\large Assignment 3: Stream Ciphers}
\author{Hugo Soares}
\date{8 de outubro de 2026}

\begin{document}
\maketitle

\section*{Questão 1}

\textbf{Ideia.} Cifrar é fazer XOR com o keystream. Como XOR é reversível,
decifrar é fazer XOR outra vez: $P = C \oplus K$.

\begin{verbatim}
def xor(a, b):
    return bytes(x ^ y for x, y in zip(a, b))

def decrypt(ciphertext, keystream):
    return xor(ciphertext, keystream).decode()
\end{verbatim}

\textbf{Pseudo-OTP.} O enunciado gera o keystream com a hora como seed:

\begin{verbatim}
random.seed(int(time.time()))
keystream = random.randbytes(n)
\end{verbatim}

Isto é mau porque o módulo \texttt{random} é determinístico: quem adivinhar
a hora reproduz o keystream todo. Com precisão de segundos, o espaço de
procura é pequeno.

\textbf{Ataque.} O ciphertext tem 37 bytes, gerado a 15/09/2026 entre as
10:00 e as 11:00 UTC. Isso dá seeds de 1789466400 a 1789470000 (3601
tentativas). Para cada seed geramos 37 bytes e fazemos XOR. Guardamos o
único resultado só com ASCII legível ($32 \le b < 127$):

\begin{verbatim}
for seed in range(start, end + 1):
    ks = random.Random(seed).randbytes(len(ct))
    pt = xor(ct, ks)
    if all(32 <= b < 127 for b in pt):
        print(seed, pt.decode())
\end{verbatim}

\begin{verbatim}
Q1 seed: 1789466511
Q1 msg: CHLORINE DOSING OVERRIDE CODE IS 7316
\end{verbatim}

\texttt{os.urandom()} evita isto porque vai buscar aleatoriedade ao sistema
operativo e não depende da hora.

\section*{Questão 2}

\textbf{Ideia.} Os comandos são conhecidos: \texttt{OPEN} (4),
\texttt{CLOSE} (5), \texttt{RESET} (5), \texttt{BYPASS} (6),
\texttt{SHUTDOWN} (8). Para cada par com igual comprimento,
$K = C \oplus P$. Mensagem curta usa só o prefixo do keystream.

\textbf{Código.} Primeiro achamos os keystreams completos com os
ciphertexts de 8 bytes (só \texttt{SHUTDOWN} tem 8). Depois cada ciphertext
é testado contra o prefixo de cada keystream:

\begin{verbatim}
keys = sorted({xor(c, p) for c in cts
               for p in cmds if len(c) == len(p) == 8})
for c in sorted(set(cts)):
    for k in keys:
        pt = xor(c, k[:len(c)])
        if pt in cmds:
            print(c.hex(), "->", pt.decode())
            break
\end{verbatim}

\textbf{Resultado.} Há 4 keystreams (120 linhas no ficheiro):

\begin{center}
\begin{tabular}{ll}
\toprule
Keystream & Hex \\
\midrule
$K_1$ & \texttt{8778e0399a6cf00a} \\
$K_2$ & \texttt{a428809b98dd0ea9} \\
$K_3$ & \texttt{d3c6299c0b2c7030} \\
$K_4$ & \texttt{fb316b0987310073} \\
\bottomrule
\end{tabular}
\end{center}

Exemplo com $K_3$: \texttt{9c966cd2} $\to$ OPEN,
\texttt{81837ad95f} $\to$ RESET, \texttt{919f79dd587f} $\to$ BYPASS,
\texttt{808e7cc84f63277e} $\to$ SHUTDOWN. Cada uma das 120 linhas bate
num único keystream. Reutilizar keystream destrói o OTP.

\section*{Questão 3}

\textbf{Ideia.} Seis mensagens com o mesmo $K$:
$C_i\oplus C_j=P_i\oplus P_j$ --- o $K$ cancela. Só há maiúsculas
($65\le b\le 90$), sem espaços, e $C_3$ tem 35 bytes.

\textbf{Testar pequenas hipóteses.} Para cada posição $i$, o byte $K[i]$
tem de dar letra em todos os 6 textos:

\begin{verbatim}
options = [[k for k in range(256)
  if all(65 <= (ct[pos] ^ k) <= 90 for ct in cts)]
  for pos in range(35)]
\end{verbatim}

Isso deixa 8 a 18 candidatos por posição, o que por si só não decide
nada. O que decide é testar palavras: se a palavra na posição $p$ de
$C_3$ implicar um $K[i]$ que não está em \texttt{options}, é rejeitada.

No início da frase, \texttt{QUE}, \texttt{NAO} e \texttt{UMA} são
aberturas comuns. \texttt{QUE} daria a $T_2$ o início \texttt{E[M} e
\texttt{NAO} daria a $T_1$ \texttt{T\_K} (o \texttt{[} e o \texttt{\_}
são bytes que não são letra); ambos violam a regra. \texttt{UMA} passa:

\begin{verbatim}
T1: OSE   T2: ACI   T3: UMA
T4: ACI   T5: CLA   T6: OSA
\end{verbatim}

É um ponto de partida útil, não a prova de que \texttt{UMA} é a única
abertura possível.

\textbf{Seguir ``uma''.} A seguir a ``uma'' testam-se nomes na posição 3.
\texttt{COISA} passa a regra das maiúsculas, mas dá \texttt{GUMAK} em
$T_1$ e \texttt{FUIAA} em $T_2$, que não encaixam nos inícios já
visíveis. \texttt{CHAVE} é natural num exercício sobre chaves
reutilizadas e dá muito melhor:

\begin{verbatim}
T1: OSEGREDO   T2: ACIFRADE   T3: UMACHAVE
T4: ACIFRADE   T5: CLAUDESH   T6: OSATAQUE
\end{verbatim}

Várias mensagens ganham palavras reconhecíveis ao mesmo tempo (``o
segredo'', ``a cifra de'', ``os ataque\ldots''). É apoio mais forte do
que uma só frase plausível em $C_3$; não é preciso adivinhar nomes nem
frases completas.

\textbf{Descrever a chave.} Na posição 8, \texttt{SECRETA} falha a regra
das letras. \texttt{FORTE} passa, mas dá \texttt{PORMM} em $T_1$, o que
produz \texttt{OSEGREDOPORMM} e não se lê naturalmente. Como o exercício
é sobre reutilização de chaves, outra descrição razoável é
\texttt{REUTILIZADA}:

\begin{verbatim}
T1: OSEGREDODEUMACIFRAD
T2: ACIFRADEVERNAMEAUNI
T3: UMACHAVEREUTILIZADA
T4: ACIFRADEVERNAMFOIPA
T5: CLAUDESHANNONDEMONS
T6: OSATAQUESDETEXTOCON
\end{verbatim}

Isto completa frases inteiras nos outros textos (``o segredo de uma
cifra'', ``os ataques de texto con\ldots''). É por isso que se fica com
\texttt{REUTILIZADA} em vez de \texttt{FORTE}. Continua a ser uma
escolha linguística, não uma consequência da regra das maiúsculas.

\textbf{Fechar a frase.} Na posição 19, \texttt{NUNCA} falha a regra e
\texttt{DEIXA} passa, estendendo $T_1$ a
\texttt{OSEGREDODEUMACIFRADEVERE}, $T_2$ a
\texttt{ACIFRADEVERNAMEAUNICAQUE} e $T_6$ a
\texttt{OSATAQUESDETEXTOCONHECID}. Depois ``uma chave reutilizada
deixa'' pede ``de ser segura'', testado em três partes: \texttt{DE} na
posição 24, \texttt{SER} na 26 e \texttt{SEGURA} na 29.

\textbf{Código.}

\begin{verbatim}
def try_word(pos, word):
    ks = [c ^ p for c, p in zip(cts[2][pos:], word.encode())]
    if any(k not in options[pos + j] for j, k in enumerate(ks)):
        return None  # implicaria um byte que nao e letra
    return [bytes(ct[pos + j] ^ k for j, k in enumerate(ks)).decode()
            for ct in cts]

trials = [(0, "QUE"), (0, "NAO"), (0, "UMA"),
          (3, "COISA"), (3, "CHAVE"),
          (8, "SECRETA"), (8, "FORTE"), (8, "REUTILIZADA"),
          (19, "NUNCA"), (19, "DEIXA"),
          (24, "DE"), (26, "SER"), (29, "SEGURA")]

accepted = [(0, "UMA"), (3, "CHAVE"), (8, "REUTILIZADA"),
            (19, "DEIXA"), (24, "DE"), (26, "SER"), (29, "SEGURA")]
keystream = [None] * 35
for pos, word in accepted:
    assert try_word(pos, word) is not None
    for j, p in enumerate(word.encode()):
        keystream[pos + j] = cts[2][pos + j] ^ p
\end{verbatim}

\texttt{options} guarda os bytes possíveis de $K$ por posição.
\texttt{try\_word} calcula os bytes implicados pela hipótese e devolve
\texttt{None} se algum falhar a regra; senão devolve o pedaço
correspondente de cada texto claro, para comparar. \texttt{accepted} são
as escolhas explicadas acima --- não há escolha automática de palavras.

\textbf{Resultado.} As partes cobrem os 35 bytes de $C_3$ (as outras
cinco continuam para lá disso):

\begin{verbatim}
T1: OSEGREDODEUMACIFRADEVERESIDIRAPENAS
T2: ACIFRADEVERNAMEAUNICAQUEOFERECESIGI
T3: UMACHAVEREUTILIZADADEIXADESERSEGURA
T4: ACIFRADEVERNAMFOIPATENTEADAEMMILNOV
T5: CLAUDESHANNONDEMONSTROUQUEOSIGILOPE
T6: OSATAQUESDETEXTOCONHECIDOEXPLORAMAR
\end{verbatim}

O texto claro completo de $C_3$ é
\texttt{UMACHAVEREUTILIZADADEIXADESERSEGURA} (``Uma chave reutilizada
deixa de ser segura''). O $K$ de 35 bytes é idêntico a partir dos 6
textos, o que confirma a decifra.

\section*{Questão 4}

\textbf{Ideia.} $C = P \oplus K$. Para trocar $P$ por $P'$ sem saber $K$,
basta anular $P$ e pôr $P'$: $C' = C \oplus P \oplus P'$.

\begin{verbatim}
c = bytes.fromhex("caffea...0071a")  # 30 bytes
p = b"PLC=07;CMD=SET;NAOH_PPM=00100;"
pp = b"PLC=07;CMD=SET;NAOH_PPM=11100;"
print(xor(c, xor(p, pp)).hex())
\end{verbatim}

\begin{verbatim}
Q4 forged: caffea70ef5b03e9f898bebf13e125a8b81e1a75aea1f20e166abbc0071a
\end{verbatim}

Só mudam os bytes 24 e 25 (\texttt{17 6b} $\to$ \texttt{16 6a}), que são os
dois primeiros dígitos de \texttt{00100} $\to$ \texttt{11100}. O ChaCha20
não foi quebrado: o problema é a falta de autenticação. Sem MAC, virar bits
no ciphertext vira os mesmos bits no plaintext. Com ChaCha20-Poly1305 a
forja seria detetada.

\section*{Questão 5}

\textbf{Ideia.} Mesmo nonce + mesma chave = mesmo keystream $K$. Como
conhecemos $P_1$ (heartbeat de 50 bytes), $K = C_1 \oplus P_1$. Depois
$P_2 = C_2 \oplus K$.

\begin{verbatim}
p1 = b"HEARTBEAT;PLC=07;STATUS=OK;NAOH_PPM=00100;PH=07.4;"
k = xor(c1, p1)
print(xor(c2, k[:len(c2)]).decode())
\end{verbatim}

\begin{verbatim}
Q5 msg: PLC=07;CMD=SETPWD;PWD=Tr3at!Pl4nt#26;
\end{verbatim}

$P_2$ tem 37 bytes e o heartbeat tem 50, por isso recupera-se tudo. Se $P_2$
fosse maior que 50 bytes, só os primeiros 50 se recuperavam; o resto
precisaria de mais plaintext conhecido.

\section*{Conclusão}

Keystream previsível, reutilizado ou sem autenticação chega para atacar,
sem quebrar o algoritmo. Usar gerador seguro, nonce único e modo
autenticado (ChaCha20-Poly1305).

\end{document}
